\scanpage{15}

$\varphi^-(x,C^-y)$ sur $V^+$ resp. $V^-$ sont définies positives, $C^+$ et $C^-$ étant des \deleted{\unsure{invol}} involutions orthogonales, on conclut que $C^+=\mathrm{id}$, $C^-=-\mathrm{id}$, donc $C$ est dans le centre de \deleted{$O(\varphi)$, \unsure{$\varphi$}} $O(V^+)\times O(V^-)$ donc dans le centre de \unsure{$G$|$G'$}.

Cas $\unsure{y|\gamma|\eta}\ne1$. Alors $gCg^{-1}\in\{C,yC\}$, donc $G'$ \deleted{\unsure{général} ($y$ pas néc. $=1$). On voit déjà que} \deleted{\unsure{est un morphisme} $G\to$} \deleted{$g\mapsto[C,g]=CgC^{-1}g^{-1}\in\{1,y\}$ pour tt $g\in G'(\mathbb R)$.}

\deleted{Soit une représentation \inserted{simple non nulle} de $G$ telle que $y$ y opère par $-\mathrm{id}$. il suffit de prouver que dans l'image, $C$ devient central. Donc on peut supposer $G\subset\operatorname{Aut}_{\mathbb R}(M)$, $y$ opérant par $-\mathrm{id}$. Donc $C$ définit sur $M$ structure complexe, et $G\subset\operatorname{Aut}_{\mathbb C}(M)$, La forme $\varphi$ de $C$-polarisation satisfait à}

opère sur $\{C,yC\}$, \deleted{\unsure{et il y a un sous-groupe}} soit $G''$ le stabilisateur de \unsure{$G$|$C$}, alors $G''$ est d'indice 1 ou 2, \deleted{\unsure{On obtient}} \inserted{et} \deleted{\unsure{on voit que l'on a}} $gCg^{-1}=\varepsilon(g)C$\quad $(\varepsilon:G\longrightarrow\mu_2)$

\[
\varphi^C(gx,gy)=\varphi(gx,Cgy)=\varepsilon(g)\,\varphi(gx,gCy)=\varepsilon(g)\varphi^C(x,y).
\]

\deleted{Donc $G'\subset\operatorname{Sp}^{\unsure{1|\prime}}(\varphi^C)$}

Si on avait $\varepsilon(g)=-1$, on aurait $\varphi^C(gx,gx)=-\varphi^C(x,x)$, et si $x\ne0$ les deux mbs seraient resp $>0$ et $<0$, absurde.

% archive sheet number: 14
